\chapter{Closed-Loop Simulation and Evaluation Methodology}
\label{ch:sim}

\section{Purpose and Design Principles}

All quantitative end-to-end claims come from a closed-loop statistical
simulator, a synthetic dynamic-graph world inspired by rover networking rather
than a wireless channel model. Three principles governed it: \emph{causality}
(delivery runs over the kept graph and execution depends on hidden reliability,
so both pruning and assignment affect success); \emph{hidden truth} (the
simulator knows each rover's reliability but the pipeline never sees it, so
learned trust can be scored against the truth); and \emph{paired comparisons}
(every arm runs on identical scenarios from the same seed).

\section{The World Model}

\subsubsection{Rovers and mobility} Rovers start at uniformly random
positions in a $10\times10$ area and take a Gaussian random-walk step (standard
deviation 0.8) every cycle, clipped to the area. Each rover has a battery
that starts between 50 and 100\%, drains by 1 to 5\% per cycle and recharges
by 3\% when lightly loaded; a load that rises by 0.15 with each assignment and
halves every cycle; and a static capability between 0.4 and 1.0.

\subsubsection{Hidden reliability} Each rover's reliability is drawn from a
base level between 0.35 and 0.75 plus a spatially correlated field made of
three Gaussian bumps of random sign and amplitude up to 0.3 (width 2.5),
clipped to $[0.05,0.98]$. Nearby rovers therefore share good or bad
conditions, which is what makes neighbourhood information useful to a GNN.
Optionally, reliability drifts as a random walk with a chosen standard
deviation per cycle, so trust must track a moving target.

\subsubsection{Tasks and mission roster} A fixed number of tasks arrives
each cycle (two by default), each with a criticality drawn between 0.2 and
0.9. A fixed mission-critical roster, the root plus a few randomly chosen
rovers, is drawn once per scenario. The roster is deliberately not redefined
every cycle as ``whoever happens to be reachable'', which would make the
connectivity guarantee trivially satisfied.

\section{The Radio Model}

Two rovers can link when they are within the communication range. The chance
that a link is heard in a given cycle is $1-\pi$, where $\pi$ is the packet
loss, set to 5\% unless stated otherwise. An observed link refreshes its quality,
$L_e=\mathrm{clip}\bigl(1-d_{ij}/d_{\max}+\mathcal N(0,0.05^2)\bigr)$, and its
timestamp; an unobserved link keeps its old quality and grows stale, with
freshness $A_e=\exp(-\Delta t_e/3)$ in cycles. The edge cost follows
Eq.~\eqref{eq:arch:cost}. So packet loss does not remove links from the
graph. What it does is make the pipeline's information about them older and
less trustworthy, which is also what tends to go wrong with real radios.

\section{Faults}

We test the trust model and the pruner with scripted faults. A
\emph{reliability drop} lowers the hidden reliability of one rover, or of all
of them. An \emph{intermittent} fault switches that drop on and off every
other cycle. A \emph{packet-loss spike} makes the link information of the
whole fleet go stale. A \emph{battery drain} empties one rover until it can no
longer carry out tasks, and a \emph{position jump} suddenly moves a rover to a
new place. In the default scenario the most reliable rover is degraded for
four cycles, and a packet-loss spike starts one cycle after that. This is the
hardest case for a system that has learned to depend on that rover.



\section{Experimental Arms}

\subsection{Pruner Arms}

Six pruners share the same utility, Eq.~\eqref{eq:arch:utility}, so they
differ only in how they choose:

\begin{itemize}
\item \textbf{full graph:} keep every candidate link;
\item \textbf{threshold:} keep links with $U_e\ge0.4$;
\item \textbf{top-$K$:} keep the best half of the links by utility;
\item \textbf{threshold+repair} and \textbf{top-$K$+repair:} the same greedy
choice followed by the same feasibility repair the QUBO path uses, enforcing
degree caps, the budget and connectivity of the mission roster; if the budget
is infeasible, the repair falls back to a minimum-cost spanning tree and
reports the residual violation rather than returning an empty graph;
\item \textbf{QUBO:} the projected annealer with the relaxation ladder of
Section~\ref{sec:sol:ladder}.
\end{itemize}

The repaired arms are essential for fairness. A table that compares the
constrained QUBO only against constraint-blind pruners would credit the QUBO
with guarantees its opponents never tried to provide.

\subsection{Budget and Degree Settings}

The budget is either a fixed number of units $K$ or a fraction $\beta_K$ of
what the full graph costs in the current cycle. The second option keeps the
pressure similar from one cycle to the next even when the graphs are
different. Unless an experiment changes it, the degree cap is three. Section~\ref{sec:res:audit} explains why the second budget mode
became necessary.

\subsection{Cloud Amortisation}

The QUBO can be solved every $k$ cycles instead of every cycle. Between
solves the last decision is reused, intersected with the live graph so that
no vanished link is ever kept. This models a cloud solver feeding a fog
device on a slower cadence.

\subsection{Assignment Arms}

The four assignment arms of Chapter~\ref{ch:gnn}, GNN, no-graph, trust-only
and random, share gates, delivery awareness and capacity, and differ only in
the scorer.

\section{Metrics}

Table~\ref{tab:sim:metrics} defines every metric reported in
Chapter~\ref{ch:results}. Four of them, offload, delivery failure and the
precision and recall of delivery-critical pruning, are usually missing from
pruning studies and were added specifically to test what the guarantees buy.

\begin{table}[htbp]
\centering
\caption{Metrics computed by the simulator.}
\label{tab:sim:metrics}
\small
\begin{tabularx}{\linewidth}{@{}l X@{}}
\toprule
Metric & Definition \\
\midrule
Mission success & fraction of tasks both executed and delivered \\
Offload rate & fraction of task slots actually assigned on the kept graph \\
Delivery-failure rate & failed deliveries divided by attempted deliveries \\
Pruning precision / recall & agreement of kept edges with \emph{delivery-critical} edges (below) \\
Trust Spearman & rank correlation between learned trust and hidden reliability \\
Trust AUC & how well trust separates rovers above and below median reliability \\
Violations & per-run counts of excess degree, budget overruns and disconnected mission rovers \\
Stage latency & median of repeated timings of build, solve, GNN and assignment \\
Recovery time & cycles until success returns to its pre-fault level \\
Relaxation usage & cycles spent on each rung of the relaxation ladder \\
\bottomrule
\end{tabularx}
\end{table}

\begin{definition}[Delivery-critical edge]
In a given cycle, a candidate edge is \emph{delivery-critical} if removing it
from the full candidate graph strictly reduces the widest-path delivery
factor $p_i$ of at least one mission-critical rover.
\end{definition}

The labels are recomputed every cycle on the live candidate graph, so they are
a counterfactual ground truth that does not depend on any pruner. Recall then
asks how many of the links that matter for delivery a pruner kept, and
precision how many of its kept links were ones that matter.

\section{Statistical Methodology}

Results are averaged over seeds, and differences between arms are always
computed \emph{per seed}, because two arms run on the same scenario share
most of their randomness. For per-seed differences $d_1,\dots,d_S$ the report
gives three numbers.

\begin{enumerate}[label=\arabic*.]
\item The mean difference $\bar d$.
\item A paired bootstrap 95\% confidence interval: resample the $S$ differences
with replacement 10\,000 times and take the 2.5th and 97.5th percentiles of
the resampled means~\cite{efron1993}.
\item A sign-flip permutation $p$-value: under the null hypothesis of no
difference, each $d_s$ is equally likely to have either sign, so the
$p$-value is the fraction of 20\,000 random sign assignments whose mean is at
least as extreme as $\bar d$.
\end{enumerate}

A difference is called significant when the interval excludes zero and
$p<0.05$. With 8 to 16 seeds per arm the power is modest, and the report
says so wherever a margin is not significant rather than drawing a trend
through noise.

\section{Scaling and Pressure Experiments}

Two further experiment drivers run outside the closed loop.

\subsubsection{Scaling} Range-limited sparse graphs, in which each rover
links to its nearest neighbours as real short-range radios would, are built
for fleets of 10 to 50 rovers. At each size the driver records the numbers of
edges and variables, the memory used by the matrix, the time to build and to
solve, the energy reached by the projected annealer and by a repaired greedy
baseline, and the GNN's latency on both the full and the pruned graph. Connectivity is enforced only up to 20 rovers,
because the dense flow encoding grows quadratically, and the report never
draws one curve across that transition.

\subsubsection{Constraint pressure} In the pressure sweep we keep the
scenarios the same and tighten one constraint at a time. The budget fraction
takes the values $\beta_K\in\{0.35,0.5,0.75,1.0\}$, the degree cap
$D\in\{2,3,4\}$ and the roster size $|R|\in\{3,5,8\}$. Every setting is run
with six pruner arms, 16 seeds and 20 cycles per run. Results are saved after
each run, so a sweep that gets interrupted can carry on from where it
stopped.

We measured all timings on an AMD Ryzen 7 5800H (16 threads, 8~GB of memory)
using Python 3.14 and NumPy with its reference linear algebra library. The
conclusions rest on the \emph{ratios} between times more than on the absolute
values.
